\documentclass[usenames,dvipsnames,french]{beamer}

\mode<presentation> {

\usetheme{Montpellier}
}

\usepackage{graphicx} % Allows including images
\usepackage[frenchb]{babel}
\usepackage[T1]{fontenc}
\usepackage[utf8]{inputenc}
% \usepackage[demo]{graphicx}
% \usepackage{caption}

\usepackage{subcaption}
% \usepackage{subfig}
% \usepackage{algorithm,algorithmic}
\usepackage{algorithm2e}
\usepackage{listings}
% \usepackage{ntheorem}
\usepackage{booktabs} % Allows the use of \toprule, \midrule and \bottomrule in tables
\setcounter{tocdepth}{2}
\graphicspath{{../Figures/}}
\DeclareMathOperator*{\argmin}{arg\,min}
\DeclareMathOperator*{\argmax}{arg\,max}
\newtheorem*{proposition}{Proposition}

% exercices comptés
\newcounter{numexos}%Création d'un compteur qui s'appelle numexos
\setcounter{numexos}{0}%initialisation du compteur
\newcommand{\exercice}[1]{%Création d'une macro ayant un paramètre
\addtocounter{numexos}{1}%chaque fois que cette macro est appelée, elle ajoute 1 au compteur numexos
\textcolor{DarkOrchid}{Exercice\,\thenumexos\,:}\,#1%Met en rouge Exercice et la valeur du compteur appelée par \thenumeexos
}
%----------------------------------------------------------------------------------------
%	TITLE PAGE
%----------------------------------------------------------------------------------------

\setbeamertemplate{title page}{
    \centering

    \Large\bfseries
    FTML: train and test set, regression.
    

\begin{figure}[!h]
\includegraphics[width=8cm]{optimal_linear_regression_standard_deviation_20_0}
\end{figure}

    }
\begin{document}

\frame{\titlepage}

\begin{frame}{Content}

    \tableofcontents

\end{frame}

\section*{Introduction}

\begin{frame}{Linear regression}

    Linear regression is one of the most elementary methods used in ML
    regression problems. It is useful for many applications, and is often a
    component of more complex methods.

    We will use is to illustrate several important aspects of ML that are also
    encountered when using other methods (kernels, trees, neural networks, etc.), namely:
    \begin{itemize}
    	\item the train and test sets.
	\item overfitting (when the dimension $d$ is not really small).
    \end{itemize}

\end{frame}

\section{Train sets and test sets, loss functions}%
\label{sec:Train sets and test sets}

\begin{frame}{Example problem}

    We want to predict the power that needs to be produced by a power plant in
    a city, as a function of the temperature only.

    \begin{figure}[htpb]
        \centering
        \includegraphics[width=0.6\textwidth]{dataset_standard_deviation_20_0}
        \caption{Dataset}
        \label{fig:dataset}
    \end{figure}
\end{frame}

\begin{frame}

    The only information we have is a collection of $n$ \textbf{samples}, called the \textbf{dataset}. Each sample consists in two values:

    \begin{itemize}
	\item the temperature in °C (the \textbf{feature}).
	\item the power consumption in W (the \textbf{label}).
    \end{itemize}

    \begin{figure}[htpb]
        \centering
        \includegraphics[width=0.4\textwidth]{dataset_standard_deviation_20_0}
        \caption{Dataset}
        \label{fig:dataset}
    \end{figure}
    \exercice{} What is $d$ in this example ?
\end{frame}

\begin{frame}{Supervised learning}

    Based on this \textbf{finite} dataset, our objective is to produce a function, noted $\tilde{f}$, called the \textbf{model}, the \textbf{predictor}, or the \textbf{estimator}, that will map a temperature to a power consumption.

    \begin{figure}[htpb]
        \centering
        \includegraphics[width=0.3\textwidth]{dataset_standard_deviation_20_0}
        \caption{Dataset}
        \label{fig:dataset}
    \end{figure}

    The problem of supervised learning is that we want this function to perform well on \textbf{new, previously unseen samples}. It is not useful in itsself to perform well on the dataset !

\end{frame}

\begin{frame}{Supervised learning}

    Based on this \textbf{finite} dataset, our objective is to produce a function, noted $\tilde{f}$, that will map a temperature to a power consumption.

    \begin{figure}[htpb]
        \centering
        \includegraphics[width=0.3\textwidth]{dataset_standard_deviation_20_0}
        \caption{Dataset}
        \label{fig:dataset}
    \end{figure}

    \exercice{} How could we \textbf{estimate} the performance of $\tilde{f}$ on unseen data, using the dataset ?
    
\end{frame}

\begin{frame}{Train set and set set}
    We split the dataset in two parts:
    \begin{itemize}
	\item the \textbf{train set} will be used to learn (train $\tilde{f}$)
	\item the \textbf{test set} will be used to estimate the performance of $\tilde{f}$ on unseen data.
    \end{itemize}
    \begin{figure}[htpb]
        \centering
        \includegraphics[width=0.5\textwidth]{dataset_standard_deviation_20_0_splitted}
        \caption{Splitted Dataset}
        \label{fig:dataset}
    \end{figure}
\end{frame}

\begin{frame}{Loss functions}

    To evaluate the performance of $\tilde{f}$, we use a \textbf{loss function}. The loss function is be a measure of the \textbf{discrepancy} between a prediction $\tilde{f}(x)$ and a \text{label} $y$ (that are both real numbers). In regression, the most common loss function is the \textbf{squared loss}:
	    \begin{equation}
		(\tilde{f}(x)-y)^2
	    \end{equation}

	 Averaging over the whole dataset, this leads to:
	 \begin{itemize}
	     \item The train error
	    \begin{equation}
		\frac{1}{n_{\text{train}}}\sum_{(x_i, y_i)\in (X_{train}\times y_{train})} (\tilde{f}(x_i)-y_i)^2
	    \end{equation}
	     \item The test error
	    \begin{equation}
		\frac{1}{n_{\text{test}}}\sum_{(x_i, y_i)\in (X_{test}\times y_{test})} (\tilde{f}(x_i)-y_i)^2
	    \end{equation}
	 \end{itemize}
\end{frame}

\begin{frame}{Empirical risk minimization}

    The most common supervised learning process is then naturally to:
    \begin{itemize}
    	\item Find $\tilde{f}$ the as a small training error.
	\item Compute the test error of $\tilde{f}$ in order to estimate its performance on unseen data.
    \end{itemize}

    \textbf{Remark}: many of these two aspects will be discussed more deeply later in the couse (e.g. notions of validation set, or optimization error)
\end{frame}

\section{Linear regression in one dimension}
\frame{\tableofcontents[currentsection]}

\begin{frame}{Linear regression}

    The remaining step is then to find $\tilde{f}$. The first step, called \textbf{model selection}, is to decide the type of function / model use. Many types of models exist, each one having potential drawbacks, for instance:
    \begin{itemize}
	\item linear functions (which we will use in this first example).
	\item polynomial functions
	\item kernels
	\item neural networks
	\item support vector machines
    \end{itemize}

    \url{https://scikit-learn.org/stable/machine_learning_map.html} 

\end{frame}


\begin{frame}{}

    \exercice{} Why are the samples not on a line (and not on a curved line either) ?

    \begin{figure}[htpb]
        \centering
        \includegraphics[width=0.6\textwidth]{dataset_standard_deviation_20_0}
        \caption{Dataset}
        \label{fig:dataset}
    \end{figure}

\end{frame}

\begin{frame}{}

    \begin{figure}[htpb]
        \centering
        \includegraphics[width=0.3\textwidth]{dataset_standard_deviation_20_0}
        \caption{Dataset}
        \label{fig:dataset}
    \end{figure}

    The power consumption does not depend \textbf{only} on the temperature,
    but also on many other variables, that we do not have access to here:
    \begin{itemize}
        \item time in the day
        \item humidity, wind
        \item period of the year (holidays or not)
        \item other variables
    \end{itemize}

\end{frame}

\begin{frame}{Linear regression}

    Formalization:
    \begin{itemize}
        \item input space (temperature in °C): $\mathcal{X} = \mathbb{R}$
        \item output space (power consumption in W): $ \mathcal{Y} = \mathbb{R}$
	\item dataset: $D=\{(x_1, y_1), \dots, (x_n, y_n), i\in[1, n]\}$.
    \end{itemize}

    \textbf{linear regression} in dimension $1$ (as $x\in \mathbb{R}$), means that our estimator is of the form:

    \begin{equation}
        \tilde{f}(x) = \theta x + b
    \end{equation}

    with $\theta\in \mathbb{R}$, $b\in \mathbb{R}$.

\end{frame}

\begin{frame}{Empirical risk minimization}

    With the squared loss, and this function $ \tilde{f}$ the \textbf{train error} is:

    \begin{equation}
	R_{n_{train}}(\theta, b) = \frac{1}{n_{train}} \sum_{(x_i, y_i)\in (X_{train}\times y_{train})} (\theta x_i+b-y_i)^2
    \end{equation}

    The \textbf{optimization problem} is to find $\theta$ and $b$ such that the empirical risk $R_{n_{train}}(\theta, b)$ has the 
    \textbf{smallest possible value}. 
\end{frame}

\begin{frame}{Analytic solutions}

    For some problems, like this one, it is possible to explicitely compute the
    optimal solution.
    \\

    Thanks to the convexity and differentiability of
    $R_n(\theta)$ with respect to $\theta$, the points optimizing the empirical risk are obtained by
    finding $ (\theta^*, b^*)$ such that the gradient cancels (more on that
    later in the course).

    \begin{equation}
        \nabla_{(\theta, b)} R_n(\theta^*, b^*) = 0
    \end{equation}
\end{frame}

\begin{frame}{Computing the optimal values}

    \exercice{}
    Compute the gradient and find the values $\theta^*$ and $b^*$ that cancel it.
\end{frame}

\begin{frame}{Overfitting}

    In the first practical session, we will observe the interplay between $n$, $d$, and the amount of overfitting.

    \begin{figure}[htpb]
    	\centering
    	\includegraphics[width=0.7\textwidth]{polyomial_regression_std_20_0_degree_15}
    	\label{fig:}
    \end{figure}

\end{frame}

\end{document} 
